\pagebreak
\section{Метод решения}
Корпус --- 1\,000\,000 документов из ЛР1 (\texttt{lab1/data\_1m/text}: 620\,164 статьи Wikipedia, 296\,250 постов Reddit, 48\,768 вопросов Travel Stack Exchange с ответами и 34\,818 статей Wikivoyage; английский язык, тематика --- туризм).

\textbf{Термины.} Текст делится на токены токенизатором из ЛР3 (\texttt{lab3/build/tokenizer}): токен --- непрерывная последовательность букв и цифр Юникода, регистр понижается. Термин --- токен после понижения регистра; стемминг и лемматизация не применяются, так что \texttt{museum} и \texttt{museums} --- разные термины.

\textbf{Частоты.} Скрипт \texttt{scripts/zipf.py} (Python~3, стандартная библиотека; запуск \texttt{zipf.py <корпус> <каталог результатов>}) подаёт все тексты корпуса на стандартный ввод одного процесса токенизатора, между файлами --- перевод строки, и считает частоту $f$ каждого термина (\texttt{collections.Counter}) по его выводу. Термины сортируются по убыванию частоты, равные частоты --- по алфавиту; ранг $r$ --- номер термина в этом списке, начиная с 1.

\textbf{Закон Ципфа} --- $f \cdot r = C$, т.\,е. $f(r) = C / r$. Константа $C$ взята равной частоте самого частого термина, $C = f(1)$, чтобы прямая закона проходила через первую точку. В логарифмических осях закон --- прямая с наклоном $-1$.

\textbf{Результаты скрипта} (каталог \texttt{results/}):
\begin{itemize}
    \item \texttt{freq.tsv} --- полный частотный список: ранг, термин, частота;
    \item \texttt{zipf.dat} --- точки графика: ранг, частота, $C/r$. Взяты первые ранги подряд, а дальше ранги через каждые 3\,\% ($r_{k+1} \approx 1{,}03\,r_k$) и последний ранг, всего 430 точек вместо 3,97 млн;
    \item \texttt{summary.tsv} --- число файлов, токенов, терминов, терминов с частотой 1, константа $C$;
    \item \texttt{ranks-table.tex} --- таблица для отчёта.
\end{itemize}
График строится в самом отчёте пакетом \texttt{pgfplots} из \texttt{zipf.dat}, обе оси логарифмические.
